\begin{abstract}
Prescriptive maintenance asks one edge controller to run heavy analytics and, at the same time, act on a machine within a hard deadline. We study where such a 10\,ms control loop should run on the Arduino UNO Q, a board that pairs a quad-core Linux processor with a Zephyr microcontroller. We compare four configurations under identical work: a Linux thread under the default fair scheduler (A), the same thread under \texttt{SCHED\_FIFO} with locked memory (B), B on a core reserved by affinity (B-iso), and the loop on the microcontroller with setpoints sent from Linux (C). Loads range from \texttt{stress-ng} to realistic edge work: bulk Wi-Fi traffic and a containerised CNN fed by a vibration stream. One lateness definition is applied to every configuration. Under the heaviest realistic load, the typical worst case per 10{,}000 cycles was 5.6\,ms for A, 0.31\,ms for B, 1.46\,ms for B-iso and 9\us{} for C; every Linux configuration missed the 1\,ms deadline (3.5\,\%, 0.025\,\% and 0.23\,\% of cycles), while C missed none in 20{,}000 cycles and delivered all 1{,}920 speed commands to an emulated drive. Two results explain the gap. \texttt{SCHED\_FIFO} preempts programs but not the kernel's interrupt and softirq work, which is why reserving a core made B worse. On the microcontroller, thread priority decided the outcome: with library defaults the communication thread delayed the loop by up to 0.8\,ms, and raising the loop above it reduced this about 100-fold. We report 95\,\% intervals for all headline figures.
\end{abstract}

\begin{IEEEkeywords}
real-time systems, Linux scheduling, SCHED\_FIFO, Zephyr RTOS, asymmetric multiprocessing, edge computing, prescriptive maintenance
\end{IEEEkeywords}
